\chapter{Execution model}

This chapter describes how threads run, block, and yield --- the runtime that
the async-syscall ABI and the sitas reactor operate on top of.

\section{Thread lifecycle}

A thread is created with \texttt{spawn\_thread(asid, entry\_point)}. The
\texttt{asid} determines whether it is a kernel thread (\texttt{asid == 0},
runs at EL1 / ring 0) or a user thread (\texttt{asid > 0}, drops to EL0 via
\texttt{user\_trampoline}).

\subsection{Entry} The trampoline unmasks interrupts, calls the entry function,
then calls \texttt{abort()} when the entry returns. \texttt{abort()} removes
the thread from the scheduler and moves it to \texttt{DEAD\_THREADS}; the
reaper drops it (freeing its stack) from the next thread's context, after the
context switch.

\subsection{Blocking} A thread blocks by calling \texttt{block\_thread(tid,
\&observable)} through a higher-level primitive (\texttt{sleep}, \texttt{wait}).
The call marks the thread Blocked, registers its \texttt{Waker} as an observer
on the event, removes it from the LP's run queue, and \texttt{yield\_lp()} saves
the thread's execution context and switches away. When the event fires, the
waker submits the thread back to the run queue; the thread resumes exactly
where it blocked.

\subsection{Waking} \texttt{Waker::notify(tid)} calls
\texttt{submit\_ready\_thread(tid)}, which adds the thread to the least-loaded
LP's run queue and sends a wakeup IPI if that LP is idle and is not the
calling LP. The thread runs at the next scheduling opportunity.

\subsection{Yielding} \texttt{yield\_lp()} sets the context-switch-pending flag
and calls \texttt{cond\_yield\_lp()}. If the flag is set and another thread is
runnable, the current thread's callee-saved registers and stack pointer are
saved to its \texttt{ThreadContext}, and the next thread's are restored.

\section{The quantum timer}

Each LP's \texttt{RoundRobin} scheduler arms a periodic quantum
\texttt{TimerEvent} (10 ms by default). When the timer fires, the observer sets
the pending flag; the next \texttt{yield\_lp()} (or the IRQ dispatcher's tail)
performs the context switch.

The quantum is re-armed after each context switch (or when the flag is cleared
with nothing to switch to). An \texttt{armed} flag ensures exactly one quantum
event is in flight at any time, preventing unbounded timer-queue growth from
manual yields. The timer follows the same block-wake pattern as any other
event: the hardware timer fires, the IRQ dispatcher processes the queue,
\texttt{LpTimer::start()} re-arms for the next deadline.

\section{The block-yield-wake cycle}

This is the fundamental loop of the async runtime:

\begin{enumerate}
    \item \textbf{Worker thread} calls \texttt{sleep(50ms)}. The timer event
        is created, the thread is blocked on it, the event is enqueued in the
        per-LP timer queue, the hardware timer is armed for the nearest
        deadline, and \texttt{yield\_lp()} switches away.
    \item \textbf{Other threads} run (or the LP idles). The timer IRQ fires
        at each quantum; the event-queue processes expired events.
    \item \textbf{50 ms elapse.} The sleep event is at the front of the timer
        queue. \texttt{process\_events} pops it, calls \texttt{signal()}, which
        upgrades the observer \texttt{Weak} and calls \texttt{Waker::notify()},
        which calls \texttt{submit\_ready\_thread(tid)}.
    \item \textbf{Worker} is now \texttt{Ready}. At the next scheduling point
        it runs, resuming immediately after its \texttt{yield\_lp()} inside
        \texttt{sleep()}.
\end{enumerate}

The block-wake cycle for \texttt{completion::wait} is identical:
\texttt{complete()} signals the completion's observers (instead of a timer
signal), waking the blocked thread.

\section{Context saving}

\texttt{switch\_ctx(curr\_sp\_ptr, next\_sp\_ptr)} is the architecture-specific
register-and-stack swap (AArch64: \texttt{x19-x30} and \texttt{TTBR0\_EL1};
x86-64: callee-saved GPRs and \texttt{CR3}). It is called with both scheduler
locks released, because it may permanently abandon the current stack (on the
initial switch away from boot context). The post-switch code reaps dead
threads.

\endinput
